% Zone „Das kennst du schon" – Prozentrechnung, Kl. 7, schwach
\einheitenkopf[zone][Kennst du schon]{Das kennst du schon}
\setcounter{aufgabe}{0}

\begin{aufgabe}{Ich kann ablesen, wie viel Prozent eines Streifens gefärbt sind.}
Lies ab, wie viel Prozent gefärbt sind.

\swa{Gefärbt: \leerfeld[\%]}{\streifen{50}{$0\,\%$}{$100\,\%$}}
\swa{Gefärbt: \leerfeld[\%]}{\streifen{30}{$0\,\%$}{$100\,\%$}}
\swa{Gefärbt: \leerfeld[\%]}{\streifen[4]{75}{$0\,\%$}{$100\,\%$}}
\swa{Gefärbt: \leerfeld[\%]}{\streifen[5]{20}{$0\,\%$}{$100\,\%$}}
\end{aufgabe}

\begin{aufgabe}{Ich kann einen Anteil in Prozent in einen Streifen einzeichnen.}
Färbe den Anteil im Streifen.

\swa{$40\,\%$}{\streifen{0}{$0\,\%$}{$100\,\%$}}
\swa{$90\,\%$}{\streifen{0}{$0\,\%$}{$100\,\%$}}
\swa{$25\,\%$}{\streifen[20]{0}{$0\,\%$}{$100\,\%$}}
\swa{$50\,\%$}{\streifenleer[0]}
\swa{$20\,\%$}{\streifen[5]{0}{$0\,\%$}{$100\,\%$}}
\end{aufgabe}

\begin{aufgabe}{Ich kann von 10 \% auf mehr Prozent und auf das Ganze schließen.}
Rechne am Streifen.

\swz{$10\,\%$ sind $2\,$kg. Wie viel sind $30\,\%$?}{\streifen{10}{$0\,\%$}{$100\,\%$}}
\swz{$10\,\%$ sind $2\,$kg. Wie viel sind $100\,\%$?}{\streifen{10}{$0\,\%$}{$100\,\%$}}
\swz{$10\,\%$ sind $1{,}5\,$l. Wie viel sind $100\,\%$?}{\streifen{10}{$0\,\%$}{$100\,\%$}}
\swz{$20\,\%$ sind $8\,$m. Wie viel sind $100\,\%$?}{\streifen{20}{$0\,\%$}{$100\,\%$}}
\end{aufgabe}

\begin{aufgabe}{Ich kann einen Anteil als Bruch schreiben und auf den Nenner 100 erweitern.}
Schreibe den Anteil als Bruch. Erweitere auf den Nenner 100.

\swz{7 von 10 Kindern}{\bruchrechteck[5]{7}{10}}
\swz{1 von 4 Stücken}{\bruchrechteck[5]{1}{4}}
\swz{9 von 20 Losen}{\bruchrechteck[5]{9}{20}}
\swz[3]{6 von 30 Äpfeln}{\bruchrechteck[6]{6}{30}}
\swz{3 Jungen und 7 Mädchen: Anteil der Jungen}{\bruchrechteck[5]{0}{10}}
\end{aufgabe}

\begin{aufgabe}{Ich kann einen Bruch als Dezimalzahl schreiben und umgekehrt.}
Schreibe als Dezimalzahl.

\swz[1]{$\dfrac{3}{10}$}{\sachtabelle{ccc}{E & z & h}{\leerzelle & \leerzelle & \leerzelle}}
\swz[1]{$\dfrac{21}{100}$}{\sachtabelle{ccc}{E & z & h}{\leerzelle & \leerzelle & \leerzelle}}
\swz{$\dfrac{2}{5}$}{\sachtabelle{ccc}{E & z & h}{\leerzelle & \leerzelle & \leerzelle}}
\swz[1]{$\dfrac{6}{100}$}{\sachtabelle{ccc}{E & z & h}{\leerzelle & \leerzelle & \leerzelle}}
\anweisung{Schreibe als Bruch mit dem Nenner 10 oder 100.}
\swz[1]{$0{,}7$}{\sachtabelle{ccc}{E & z & h}{\makebox[1.2cm]{0} & \makebox[1.2cm]{7} & \leerzelle}}
\swz[1]{$0{,}35$}{\sachtabelle{ccc}{E & z & h}{\makebox[1.2cm]{0} & \makebox[1.2cm]{3} & \makebox[1.2cm]{5}}}
\end{aufgabe}

\begin{aufgabe}{Ich kann auf eine Stelle nach dem Komma runden.}
Runde auf eine Stelle nach dem Komma.

\swz[1]{$4{,}32$}{\zahlenstrahl[xmin=4.3,xmax=4.4,xstep=0.01,karo=0.8]{4.32/}}
\swz[1]{$6{,}78$}{\zahlenstrahl[xmin=6.7,xmax=6.8,xstep=0.01,karo=0.8]{6.78/}}
\swz[1]{$12{,}45$}{\zahlenstrahl[xmin=12.4,xmax=12.5,xstep=0.01,karo=0.8]{12.45/}}
\swz[1]{$3{,}96$}{\zahlenstrahl[xmin=3.9,xmax=4.0,xstep=0.01,karo=0.8]{3.96/}}
\swz[1]{Anzeige: $27{,}272727$}{\zahlenstrahl[xmin=27.2,xmax=27.3,xstep=0.01,karo=0.8]{27.2727/}}
\end{aufgabe}

\begin{aufgabe}{Ich kann durch 100 teilen und mit einer Dezimalzahl malnehmen.}
Rechne im Kopf.

\swz[1]{$500 : 100$}{\sachtabelle{ccccc}{H & Z & E & z & h}{\leerzelle & \leerzelle & \leerzelle & \leerzelle & \leerzelle}}
\swz[1]{$64 : 100$}{\sachtabelle{ccccc}{H & Z & E & z & h}{\leerzelle & \leerzelle & \leerzelle & \leerzelle & \leerzelle}}
\swz[1]{$9 : 100$}{\sachtabelle{ccccc}{H & Z & E & z & h}{\leerzelle & \leerzelle & \leerzelle & \leerzelle & \leerzelle}}
\swz[1]{$0{,}5 \cdot 30$}{$0{,}5 \cdot 30 = 5 \cdot 30 : 10$}
\swz[1]{$0{,}2 \cdot 45$}{$0{,}2 \cdot 45 = 2 \cdot 45 : 10$}
\swz[1]{$0{,}03 \cdot 200$}{$0{,}03 \cdot 200 = 3 \cdot 200 : 100$}
\end{aufgabe}

\begin{aufgabe}{Ich kann mit einer Tabelle über 1 Stück rechnen.}
Rechne mit dem Dreisatz.

\swz[1]{5 Stifte kosten $10\,€$. Wie viel kosten 3 Stifte?}{\begin{dreisatz}{Stifte}{Preis}
\dsz{5}{10\,€}
\dsp{:5}{:5}
\dsz{1}{\dsleer}
\dsp{\cdot 3}{\cdot 3}
\dsz{3}{\dsleer}
\end{dreisatz}}
\swz[1]{2 kg Äpfel kosten $6\,€$. Wie viel kosten 5 kg?}{\begin{dreisatz}{kg}{Preis}
\dsz{2}{6\,€}
\dsp{:2}{:2}
\dsz{1}{\dsleer}
\dsp{\cdot 5}{\cdot 5}
\dsz{5}{\dsleer}
\end{dreisatz}}
\swz[1]{4 m Band kosten $4{,}80\,€$. Wie viel kosten 3 m?}{\begin{dreisatz}{m}{Preis}
\dsz{4}{4{,}80\,€}
\dsp{\phantom{:0}}{\phantom{:0}}
\dsz{1}{\dsleer}
\dsp{\phantom{:0}}{\phantom{:0}}
\dsz{3}{\dsleer}
\end{dreisatz}}
\swz[1]{3 Tüten wiegen $600\,$g. Wie viel wiegen 10 Tüten?}{\begin{dreisatz}{Tüten}{Gewicht}
\dsz{3}{600\,\text{g}}
\dsp{\phantom{:0}}{\phantom{:0}}
\dsz{1}{\dsleer}
\dsp{\phantom{:0}}{\phantom{:0}}
\dsz{10}{\dsleer}
\end{dreisatz}}
\end{aufgabe}

\begin{aufgabe}{Ich kann einen Bruchteil einer Größe ausrechnen.}
Rechne.

\swz{$\frac{1}{2}$ von $30\,€$}{\streifen[2]{50}{$0\,€$}{$30\,€$}}
\swz{$\frac{1}{4}$ von $20\,$kg}{\streifen[4]{25}{$0\,$kg}{$20\,$kg}}
\swz{$\frac{2}{5}$ von $40\,$m}{\streifen[5]{40}{$0\,$m}{$40\,$m}}
\swz{$\frac{3}{10}$ von $70\,€$}{\streifen{30}{$0\,€$}{$70\,€$}}
\end{aufgabe}

\begin{aufgabe}{Ich finde den Fehler beim Umwandeln in eine Dezimalzahl.}
Mia hat umgewandelt. Finde den Fehler, benenne ihn und schreibe richtig.

\swz{Mias Rechnung:}{\rechnung{\tfrac{7}{100} &= 0{,}7}}
\swz{Mias Rechnung:}{\rechnung{0{,}05 &= \tfrac{5}{10}}}
\end{aufgabe}

\begin{aufgabe}{Ich kann Hundertstel als Dezimalzahl schreiben und umgekehrt.}
Schreibe als Dezimalzahl.

\swz[1]{$\dfrac{4}{100}$}{\sachtabelle{ccc}{E & z & h}{\leerzelle & \leerzelle & \leerzelle}}
\swz[1]{$\dfrac{40}{100}$}{\sachtabelle{ccc}{E & z & h}{\leerzelle & \leerzelle & \leerzelle}}
\anweisung{Schreibe als Bruch mit dem Nenner 100.}
\swz[1]{$0{,}08$}{\sachtabelle{ccc}{E & z & h}{\makebox[1.2cm]{0} & \makebox[1.2cm]{0} & \makebox[1.2cm]{8}}}
\swz[1]{$0{,}8$}{\sachtabelle{ccc}{E & z & h}{\makebox[1.2cm]{0} & \makebox[1.2cm]{8} & \leerzelle}}
\end{aufgabe}
